\chapter*{Abstract}
\addcontentsline{toc}{chapter}{Abstract}

\selectlanguage{english}

This thesis presents a two-year apprenticeship carried out within the Research \& Development department of Souchier-Boullet, an industrial company specialized in fire-rated compartmentation, natural smoke extraction and, for some product lines, solutions combining fire safety and acoustic performance. The work was structured around two complementary missions: improving and maintaining parametric 3D models in Autodesk Inventor, and designing SQL-based product configurators supported by an Infor Design Studio interface.

The first mission focused on engineering changes coming from production and design office feedback, with the objective of keeping digital models, drawings and bills of materials aligned with manufacturing reality. The second mission aimed at replacing a largely manual configuration workflow based on fire test reports and spreadsheets with a relational database architecture able to formalize business rules, automate calculations and secure the selection of valid product variants.

The methodology first developed during the A4 year on the Coveraccess range was consolidated through an initial extension to DOORSLIDE, then further developed during the A5 year with DOORSTEEL, PHONI LUX-AIR-PASS, ONYX and MULTIDOOR doors, and ScreenPass and CURTEX V2 compartmentation curtains. This progression confirmed the relevance of a generic data model while also revealing the diversity of technical constraints involved, from geometry and regulatory compliance to mechanical and acoustic performance. A major outcome of the second year was the dynamic integration between the SQL configurator and Autodesk Inventor, allowing the automatic opening of the appropriate 3D model according to the parameters selected in the commercial interface.

The thesis shows that these developments generated tangible benefits: better technical knowledge capitalization, lower risk of manual error, improved cross-functional communication and a stronger digital continuity between configuration and design activities. It also highlights the limitations of an integration strategy based on implicit naming conventions. The scientific chapter addresses this issue by proposing a formal interface contract between the SQL configurator and the CAD modeler, together with a validation mechanism intended to detect compatibility breaks before execution.

\paragraph{Keywords:} SQL, product configurator, CAD, Autodesk Inventor, fire compartmentation, smoke control, interoperability, industrial digital transformation.

\selectlanguage{french}
